\documentclass{paper}

\usepackage{amsmath,lstautogobble,amssymb,amsfonts, listings, fancyhdr, stmaryrd, array, tikz}
\usepackage[many]{tcolorbox}
\usepackage[T1]{fontenc}
\usepackage[utf8]{inputenc}
\usepackage[a4paper, total={170mm,257mm}, left=20mm, top=20mm]{geometry}
\usepackage{adjustbox}
\definecolor{color0}{HTML}{000000}
\definecolor{color1}{HTML}{000000}
\definecolor{bgColor}{HTML}{ffffff}
\usetikzlibrary {arrows.meta,bending} 

\newtheorem{proposition}{Proposition}
\newtheorem{definition}{Definition}
\newtheorem{preuve}{Preuve}

\definecolor{codegreen}{rgb}{0,0.6,0}
\definecolor{codegray}{rgb}{0.5,0.5,0.5}
\definecolor{codepurple}{rgb}{0.58,0,0.82}
\definecolor{backcolour}{rgb}{0.95,0.95,0.92}
\lstdefinestyle{mystyle}{
    language=caml,
    backgroundcolor=\color{backcolour},   
    commentstyle=\color{codegreen},
    keywordstyle=\color{magenta},
    numberstyle=\tiny\color{codegray},
    stringstyle=\color{codepurple},
    basicstyle=\ttfamily\footnotesize,
    breakatwhitespace=false,         
    breaklines=true,                 
    captionpos=b,                    
    keepspaces=true,                 
    numbers=left,                    
    numbersep=5pt,                  
    showspaces=false,                
    showstringspaces=false,
    showtabs=false,                  
    tabsize=2,
    autogobble=true
}

\lstset{style=mystyle}

\newcolumntype{C}{>$c<$}
\tcolorboxenvironment{proposition}{enhanced, borderline={0.8pt}{0pt}{blue}, borderline={0.4pt}{2pt}{cyan}, boxrule=0.4pt, colback=white, coltitle=black, sharp corners}
\tcolorboxenvironment{definition}{ enhanced, borderline={0.8pt}{0pt}{red}, borderline={0.4pt}{2pt}{orange}, boxrule=0.4pt, colback=white, coltitle=black, sharp corners}
\tcolorboxenvironment{preuve}{ enhanced, borderline={0.8pt}{0pt}{green}, borderline={0.4pt}{2pt}{lime}, boxrule=0.4pt, colback=white, coltitle=black, sharp corners}

\pagestyle{fancy}
\fancyhead[C]{TIPE-ENS Rapport - Dorian GIL}

\begin{document}
\setlength{\headheight}{13.07225pt}
\addtolength{\topmargin}{-1.07225pt}

\section*{Etude d'un algorithme de vérification de modèle en Logique Linéaire Temporelle}

Dans un monde dont la dépendance à divers technologies est très importante, il est important de s'assurer du bon fonctionnement de ces technologies.
C'est ainsi que la vérification de modèle rentre en jeu pour répondre à ces problèmatiques. Ce sont des techniques automatiques qui étant donné un 
modèle à nombre d'état fini d'un système et des propriétés formelles, vérifie systématiquement si ces propriétés restent vrai pour ce modèle.
Nous allons donc étudier un algorithme proposé par Rob Gerth, Doron Peled, Moshe Ya'akov Vardi, Pierre Wolper en 1995.

La bible: Simple On-the-fly Automatic Verification of Linear Temporal Logic

\section{Préliminaires}
    Nous allons pour l'instant introduire des concepts hors programme MPI afin de mener à bien notre étude.

\subsection{Logique Linéaire Temporelle}
    On modélise les différentes contraintes auxquels notre système doivent obéir par des formules logiques "dépendant du temps" d'où:
    \begin{definition}[Formule de la LTL]
        On définit les opérateurs suivants:
            \begin{itemize}
                \item $X\phi$ : $\phi$ doit être satisfaite dans l'état suivant (neXt)
                \item $\psi U\phi$ : $\psi$ doit être satisfaite dans tous les états jusqu'à un état où $\phi$ est satisfait (Until)
            \end{itemize}
        On définit $F$ l'ensemble des formules de la LTL inductivement par:
        \begin{itemize}
            \item $p\in AP \implies p\in F$
            \item si $\psi$ et $\phi$ sont des formules de LTL alors $\lnot\psi, \phi\lor\psi, \textbf{X}\psi, \phi \textbf{U}\psi$ sont des formules de LTL
        \end{itemize}
        $AP$ un ensemble fini de variables propositionnelles.
    \end{definition}

    Soit $v : \omega\times F \rightarrow \{T,F\}$ une fonction de valuation.
    
    \begin{definition}[Monde]
        On définit un tel objet comme $\omega := w_0,w_1,\dots$ une suite infinie d'état.
        On écrira $\omega^i:=w_i,w_{i+1},\dots$ un suffixe de $\omega$.
        On définit $\omega \models f$ via:
        \begin{itemize}
            \item $\omega \models a \Leftrightarrow v(w_0,a) = T$ si $a$ atomique 
            \item $\omega \models \lnot f$ si $\lnot(\omega \models f)$ 
            \item $\omega \models f\lor g$ si $\omega \models f$ ou $\omega \models g$
            \item $\omega \models \textbf{X}f$ si $\omega^1 \models f$
            \item $\omega \models f\textbf{U}g$ si $\omega \models g$ ou $\omega \models f\land\textbf{X}(f\textbf{U}g)$
        \end{itemize}
    \end{definition}

\subsection{Automate de Buchi}
    Ces formules seront transformés en automate acceptant des mots infinis, donc on doit définir de nouveaux automates:
    \begin{definition}[Automate de Buchi (BA)]
        Un automate de Büchi est un 5-uplets $(S, I, T, F, \Sigma)$ tel que
        \begin{itemize}
            \item $S$ est un ensemble fini d'état
            \item $I\subseteq S$ un ensemble d'état initiaux
            \item $F\subseteq S$ un ensemble d'état finaux
            \item $\Sigma$ Un ensemble fini de symboles appelé alphabet
            \item $T: \{S\times \Sigma\}\mapsto \mathcal{P}(S)$ Une fonction de transition.
        \end{itemize}
        Cet automate particulier accepte des séquences infinis (donc des mots infinis) ssi le mot passe par un état acceptant une infinité de fois.
    \end{definition}
    On doit pouvoir faire deux opérations sur nos automates, un test pour savoir si un automate est vide et intersecter des automates, tout en intersectant les langages naturellement.
    \begin{proposition}[Test BA Vide]
        Un BA est non vide ssi il existe un état final qui est atteignable depuis un état initial et appartient à un cycle.
        
        \textbf{Algorithme:} On considère notre automate comme un graphe orienté, qu'on décompose en CFC, puis on fait un DFS depuis les etats initiaux et on trouve une CFC non trivial tq il contient un etat final et est atteignable par un état initial.
    \end{proposition}
    \begin{preuve}
    
    \end{preuve}
    Notre algorithme à un moment un automate de Büchi généralisé, on va devoir donc le transformer en automate de Büchi usuel.
    \begin{definition}[Automate de Büchi généralisé (GBA)]
        C'est un automate de Büchi où on change l'ensemble des états finaux $F$ par $F_g\subseteq \mathcal{P}(S)$ appelé \textbf{condition d'acceptation}.
    \end{definition}
    Un GBA accepte un mot ssi il est passé une infinité de fois par un état dans chaque ensemble d'état acceptant.

    \begin{proposition}[GBA vers BA]
        Un ensemble multiple d'état acceptant peut être traduit en un ensemble en construisant un automate par une "counting construction". 
        C'est à dire si $A = (S,I, \{F_1,\dots,F_n\}, \Sigma, T)$, alors il est equivalent à $A' = (S',I', F, \Sigma, T')$ telle que
        \begin{itemize}
            \item $Q' = Q\times\{1,...,n\}$
            \item $I' = I\times\{1\}$
            \item $F'=F_1\times \{1\}$
            \item $\Delta' = \{ ( (q,i), a, (q',j) ) | (q,a,q') \in \Delta$ et si $q \in F_i$, $j= i+1 \mod n$ sinon $j=i \}$
        \end{itemize}
        qui est bien un automate de Büchi (non généralisé) qui accepte les mêmes mots.
    \end{proposition}
    \begin{preuve}
        
    \end{preuve}

\begin{proposition}[Intersection de BA]
    On définit l'automate reconnaissant l'intersection de deux langages comme $A'= (S', I', F', \Sigma, T')$ telle que
    \begin{itemize}
        \item $S'=S_1\times S_2\times \{1,2\}$
        \item $T'=T_1'\cup T_2'$
        \begin{itemize}
            \item $T_1=\{((q_1,q_2,1),a,(q_1',q_2',i)), (q_1,a,q_1')\in T_1, (q_2,a,q_2')\in T_2, q_1\in F_1 \Leftrightarrow i= 2, i\in[|1,2|]\}$
            \item $T_2=\{((q_1,q_2,2),a,(q_1',q_2',i)), (q_1,a,q_1')\in T_1, (q_2,a,q_2')\in T_2,  q_2\in F_2 \Leftrightarrow i= 1, i\in[|1,2|]\}$
        \end{itemize}
        \item $I'= I_1\times I_2 \times \{1\}$
        \item $F'= \{(q_1,q_2,2), q_2\in F_2 \}$
    \end{itemize}
    Il reconnait l'intersection des langages reconnues par les automates de Büchi
\end{proposition}
\begin{preuve}
    
\end{preuve}

\subsection{La vérification de modèle}
    \begin{definition}
        Techniques automatiques qui étant donné un modèle à nombre d'état fini d'un système et des propriétés formelles, vérifie systématiquement si ces propriétés restent vrai pour ce modèle.
    \end{definition}

    On présente ici le déroulé de l'algorithme de vérification de modèle dans notre cas :
    \begin{enumerate}
        \item Création d'un automate $A_m$ correspondant au cahier des charges d'une technologies, sous forme des différents états dans lequel la technologie sera.
        \item Création d'un automate $A_{\lnot\varphi}$ où $\varphi$ représente les differentes contraintes que notre cahier des charges doit respecter, sous forme
        d'une formule logique en logique linéaire temporelle. 
        \item Création de l'automate Syncronisé $A_m\times A_{\lnot\varphi}$
        \item Si l'automate resultant est vide, alors c'est bon, sinon on renvoie un contre exemple.
    \end{enumerate}

    Interessons nous donc à la construction de $A_m$: Intuitivement, Le langage des $A_m$ représente toutes les exécutions de $M$ et celui de $A_{\lnot\varphi}$ représente les exécutions ne satisfaisant pas la 
    formule $\varphi$. Le langage du produit est donc vide ssi tous les chemins de $M$ satisfont $\varphi$. Sinon il existe un mot en commun, qui nous donne l'erreur.

    En fait,  l'automate est donné après transformation d'un système de transition $\mathcal{M}$ en un automate $A_{\mathcal{M}}$
\section{Résolution}

\subsection{Création de $A_m$}


\subsection{Creation de $A_{\lnot\varphi}$}
    On présente ici l'algorithme de Gerth

\subsection{Création de l'automate syncronisé $A_m\times A_{\lnot\varphi}$}

\subsection{Conclusion et Commentaires}



\section{Correction}

\begin{proposition}[Correction de la création de $A_{\lnot\varphi}$]
    L'algorithme de Gerth-Valdi-Peled-Wolper est correct
\end{proposition}

\begin{preuve}
    
\end{preuve}

\begin{proposition}[Correction de l'algorithme]
    Notre processus entier de vérification de modèle est correct.
\end{proposition}

\begin{preuve}

\end{preuve}

\section{S'il y a pas assez de matière: Améliorations possibles}

\section{Bibliographie}
\begin{enumerate}
    \item Rob Gerth, Doron Peled, Moshe Y. Vardi, Pierre Wolper ; Simple On-the-fly Automatic Verification of Linear Temporal Logic ; https://orbi.uliege.be/bitstream/2268/116644/1/GPVW95-pstv.pdf
    \item Christel Baier, Joost-Pieter Katoen ; Principles of Model Checking ; The MIT Press - ISBN: 9780262026499
\end{enumerate}

\end{document}